\section{Regression}

The Python part of this exercise is available in \myhref{https://github.com/LosDelDGIIM/LosDelDGIIM.github.io/blob/main/subjects/Erasmus-DUE/GKI/Exercises/Exercise7.ipynb}{this Jupyter Notebook}.\\

\begin{ejercicio}
    What is the primary goal of linear regression?

    Linear regression is a statistical method used to model the relationship between a dependent variable and one or more independent variables. The primary goal of linear regression is to find the best-fitting line (or hyperplane in higher dimensions) that describes the relationship between the independent variables and the dependent variable. This line can then be used to make predictions about the dependent variable based on new values of the independent variables.
\end{ejercicio}

\begin{ejercicio}
    Explain the difference between independent and dependent variables in a regression model.

    In a regression model, the independent variable(s) (also known as predictor(s) or feature(s)) are the variable(s) that are used to predict the value of the dependent variable. The dependent variable (also known as the response variable or target variable) is the variable that we are trying to predict or explain. The independent variables are assumed to have an influence on the dependent variable, and the goal of regression is to quantify this relationship.
\end{ejercicio}

\begin{ejercicio}
    What are the assumptions of linear regression, and why are they important?

    The assumptions of linear regression are:
    \begin{itemize}
        \item Linearity: The relationship between the independent and dependent variables is linear.
        \item Independence: The residuals (errors) are independent of each other.
        \item No multicollinearity: The independent variables are not highly correlated with each other.
        \item Normality: The residuals are normally distributed.
        \item Homoscedasticity: The variance of the residuals is constant across all levels of the independent variables.
    \end{itemize}
    These assumptions are important because they ensure that the estimates of the regression coefficients are unbiased and efficient. Violations of these assumptions can lead to incorrect conclusions and predictions from the regression model.
\end{ejercicio}

\begin{ejercicio}
    How do Mean Squared Error (MSE) and R-squared (R2) measure the performance of a regression model?

    Mean Squared Error (MSE) is a measure of the average squared difference between the observed actual outcomes and the predicted values from the regression model. A lower MSE indicates a better fit of the model to the data.

    R-squared (R2) is a statistical measure that represents the proportion of the variance in the dependent variable that is explained by the independent variables in the model. An R2 value of 1 indicates that the model perfectly explains the variance in the dependent variable, while an R2 value of 0 indicates that the model does not explain any of the variance. Higher R2 values indicate a better fit of the model to the data.
\end{ejercicio}

\begin{ejercicio}
    What is the purpose of regularization (L1 and L2) in regression, and how do they differ?

    Regularization is a technique used to prevent overfitting in regression models by adding a penalty term to the loss function. The two most common types of regularization are L1 (Lasso) and L2 (Ridge) regularization.
    \begin{itemize}
        \item L1 regularization (Lasso) adds a penalty equal to the absolute value of the coefficients, which can lead to some coefficients being exactly zero. This means that L1 regularization can perform feature selection by effectively removing some features from the model.
        \item  L2 regularization (Ridge) adds a penalty equal to the square of the coefficients, which tends to shrink the coefficients towards zero but does not set them exactly to zero. This means that L2 regularization can help reduce the impact of less important features without completely removing them from the model, reducing the impact of multicollinearity and improving the stability of the model.
    \end{itemize}

    In summary, both L1 and L2 regularization aim to improve the generalization of the regression model by preventing overfitting, but they do so in different ways: L1 can lead to sparse models with feature selection, while L2 shrinks coefficients without eliminating features.
\end{ejercicio}